% META: {"id": "1SPE-SECDEG-CR-014", "chapitre": "1SPE-SECOND-DEGRE", "type_objet": "cours", "capacites_codes": ["C5"], "sources_inspiration": [], "status": "generated"}

\section{Inéquations du second degré}

% ============================================================
% STRATE 1 — Méthode de résolution
% ============================================================

\propriete[4 — Méthode de résolution d'une inéquation du second degré]{
Pour résoudre une inéquation du type $ax^2 + bx + c > 0$ (ou $\geq 0$, $< 0$, $\leq 0$) :

\begin{enumerate}
  \item \textbf{Ramener à zéro :} si l'inéquation n'est pas de la forme $[\ldots] \bowtie 0$, transposer tous les termes du même côté.
  \item \textbf{Calculer le discriminant} $\Delta = b^2 - 4ac$ et les racines éventuelles $x_1 \leq x_2$.
  \item \textbf{Dresser le tableau de signes} du trinôme (voir C4).
  \item \textbf{Lire la solution} sur le tableau de signes.
\end{enumerate}
}

\exemple{
Résoudre $x^2 - x - 6 > 0$.

\textbf{Étape 1.} L'inéquation est déjà ramenée à $0$.

\textbf{Étape 2.} $\Delta = 1 + 24 = 25 > 0$ ; $x_1 = \dfrac{1 - 5}{2} = -2$ ; $x_2 = \dfrac{1 + 5}{2} = 3$.

\textbf{Étape 3.} $a = 1 > 0$ :

\begin{center}
\begin{tabular}{|c|ccccccc|}
\hline
$x$ & $-\infty$ & & $-2$ & & $3$ & & $+\infty$ \\
\hline
$x^2-x-6$ & & $+$ & $0$ & $-$ & $0$ & $+$ & \\
\hline
\end{tabular}
\end{center}

\textbf{Étape 4.} $x^2 - x - 6 > 0$ sur $\boxed{]-\infty\,;\,-2[ \cup ]3\,;\,+\infty[}$.
}

\exemple{
Résoudre $2x^2 + 3x - 2 \leq 0$.

$\Delta = 9 + 16 = 25$ ; $x_1 = \dfrac{-3-5}{4} = -2$ ; $x_2 = \dfrac{-3+5}{4} = \dfrac{1}{2}$.

$a = 2 > 0$ : le trinôme est négatif ou nul entre les racines.

Solution : $\boxed{\left[-2\,;\,\dfrac{1}{2}\right]}$.
}

\exemple{
Résoudre $3x + 5 \leq x^2 + 1$.

\textbf{Étape 1 (ramener à zéro) :} $0 \leq x^2 - 3x - 4$, soit $x^2 - 3x - 4 \geq 0$.

$\Delta = 9 + 16 = 25$ ; $x_1 = -1$ ; $x_2 = 4$.

$a = 1 > 0$ : positif ou nul hors de $]-1\,;\,4[$.

Solution : $\boxed{]-\infty\,;\,-1] \cup [4\,;\,+\infty[}$.
}

\margeAppui{%
\textbf{Les 4 étapes :}\newline
1. Ramener à zéro.\newline
2. Discriminant et racines.\newline
3. Tableau de signes.\newline
4. Lire la solution.%
}

\margeAppui{%
\textbf{Rappel :} si $\Delta < 0$, le trinôme est du signe de $a$ partout : la solution est soit $\mathbb{R}$ entier, soit $\emptyset$, selon l'inégalité demandée.%
}

% ============================================================
% STRATE 2 — Erreurs fréquentes
% ============================================================

\erreurFrequente{
\textbf{Oublier de ramener à zéro avant d'étudier le signe.}

Face à $x^2 - 3 > 2x$, certains élèves étudient le signe de $x^2 - 3$ d'un côté et celui de $2x$ de l'autre, puis comparent. Cette démarche est fausse. Il faut d'abord écrire $x^2 - 2x - 3 > 0$ pour disposer d'un seul trinôme dont on étudie le signe.
}

\erreurFrequente{
\textbf{Erreur sur les bornes : inclure ou exclure les racines.}

Si l'inéquation est stricte ($>$ ou $<$), les racines ne font pas partie de l'ensemble-solution : on utilise des crochets ouverts $]$ et $[$. Si l'inéquation est large ($\geq$ ou $\leq$), les racines sont incluses : on utilise des crochets fermés $[$ et $]$.
}

% ============================================================
% STRATE 3 — Équations bicarrées
% ============================================================

\approfondissement{
\textbf{Équations bicarrées.}

Une \textbf{équation bicarrée} est de la forme $ax^4 + bx^2 + c = 0$ ($a \neq 0$). Le changement de variable $u = x^2$ ($u \geq 0$) la transforme en équation du second degré en $u$ :
\[
  au^2 + bu + c = 0.
\]
On résout cette équation en $u$, puis on revient à $x$ en résolvant $x^2 = u$ (valable seulement si $u \geq 0$).

\medskip
\textbf{Exemple.} Résoudre $x^4 - 5x^2 + 4 = 0$.

On pose $u = x^2$ : $u^2 - 5u + 4 = 0$.

$\Delta = 25 - 16 = 9$ ; $u_1 = 1$ ; $u_2 = 4$.

\begin{itemize}
  \item $u_1 = 1$ : $x^2 = 1 \Rightarrow x = \pm 1$.
  \item $u_2 = 4$ : $x^2 = 4 \Rightarrow x = \pm 2$.
\end{itemize}

L'ensemble des solutions est $\{-2\,;\,-1\,;\,1\,;\,2\}$.

\medskip
\textbf{Piège fréquent.} Si l'équation en $u$ donne une solution $u < 0$, on la rejette car $x^2 \geq 0$.
}
